\begin{proof}[Proof of Theorem~\ref{thm:pairs}]
Lemma~\ref{lem:local-control} gives
\[
 F(4u+r)\equiv D(r)\pmod{16}
 \qquad(u\in\Ztwo,\ 0\le r\le3).
\]
The four residues are $1,3,13,15$, respectively. If $k=0$, the
indices have different parity, and their values differ by $2$ modulo
$4$. If $k=1$, they have the same parity and different residues modulo
$4$, and their values differ by $4$ modulo $8$. These observations give
the first branch of the formula.

Suppose $k\ge2$. Then $M,N$ have the same residue modulo $4$.
First suppose $t_M=t_N=t$. The value $t$ is finite: the only
nonnegative integer with infinite $t_X$ is $X=0$, and the indices are
distinct. If the indices are even, their quotients by $2^t$ are odd;
if they are odd, the same is true of $(M+1)/2^t$ and $(N+1)/2^t$.
Consequently $k\ge t+1$. After interchanging the indices, assume
$M>N$, and write
\[
 h=M-N=U2^k,\qquad B=2N-M=N-h,
\]
where $U$ is positive and odd. Since $k>t$, subtracting $h$ preserves
the valuation of the even member of $N,N+1$. Thus
$v_2(B(B+1))=t$, in particular $B\ne0,-1$. Moreover,
\[
 B+U2^k=N\ge0,\qquad B+U2^{k+1}=M.
\]
Theorem~\ref{thm:main}, in its range $k\ge\max\{2,t\}$, now gives
$v_2(D(M)-D(N))=k+2t$.

If $t_M\ne t_N$, put $a=\min(t_M,t_N)$. The common residue modulo
$4$ must be $0$ or $3$, since in the other two classes both $t$-values
are $1$. Hence $a\ge2$. The two even quantities, either $M,N$ or
$M+1,N+1$, have unequal valuations, so $k=a$. In the even case,
Proposition~\ref{prop:basepoints} gives
\[
 v_2(D(M)-1)=3t_M,\qquad v_2(D(N)-1)=3t_N,
\]
with $v_2(D(0)-1)=+\infty$ when one index is zero. Subtraction of
terms of unequal valuations gives $3a$. In the odd case the same
argument uses $D(M)+1$ and $D(N)+1$. Thus in both cases the valuation
is $3a=k+2\min(t_M,t_N)$, as required.
\end{proof}

\section{The interpolating function and its image}\label{sec:image}
Write $\Ztwo^\times$ for the odd $2$-adic integers. Let $m$ denote
additive Haar measure on $\Ztwo$, normalized by $m(\Ztwo)=1$.

\begin{theorem}\label{thm:image}
The function $F:\Ztwo\longrightarrow\Ztwo$ is injective. Its only
critical points are $0$ and $-1$, and
\[
 v_2(F'(x))=2v_2(x(x+1))\qquad(x\in\Ztwo\setminus\{0,-1\}).
\]
The zeros of $F(x)-1$ at $0$ and of $F(x)+1$ at $-1$ have order
three, with unit leading Taylor coefficients. The image is the
disjoint union
\begin{align}
 F(\Ztwo)={}&\{1,-1\}\ \cup\ (3+16\Ztwo)\ \cup\ (13+16\Ztwo)
 \notag\\
 &{}\cup\bigcup_{a\ge2}
 \left[(1+2^{3a}\Ztwo^\times)\ \cup\
       (-1+2^{3a}\Ztwo^\times)\right].\label{eq:exact-image}
\end{align}
In particular, $m(F(\Ztwo))=1/7$.
\end{theorem}
\begin{proof}
First extend Theorem~\ref{thm:pairs} to distinct $x,y\in\Ztwo$.
Set $k=v_2(x-y)$ and $t_z=v_2(z(z+1))$, with
$t_0=t_{-1}=+\infty$. Choose nonnegative integers $M_j\to x$ and
$N_j\to y$. For sufficiently large $j$, $v_2(M_j-N_j)=k$. If
$k\ge2$, at least one of $t_x,t_y$ is finite, since the distinct
roots $0,-1$ differ by an odd number. The minimum of the two
$t$-values therefore also stabilizes: a finite value is locally
constant, and the $t$-value of an approximating sequence tending to
a root tends to infinity. The integer formula and continuity of $F$
give
\[
 v_2(F(x)-F(y))=
 \begin{cases}
 k+1,& k\le1,\\
 k+2\min(t_x,t_y),& k\ge2.
 \end{cases}\label{eq:padic-pairs}\tag{\theequation}
\]
Indeed, a convergent sequence all of whose terms lie in
$2^e\Ztwo^\times$ has its limit in that same closed set. The
right-hand side is finite, so $F$ is injective.

Lemma~\ref{lem:local-control} gives local analyticity on all four
residue classes. For $x\ne0,-1$, let $t=v_2(x(x+1))$. Whenever
$v_2(h)>t$ and $v_2(h)\ge2$, one has $t_{x+h}=t$. Applying
\eqref{eq:padic-pairs} and dividing by $h$ yields
\[
 v_2\left(\frac{F(x+h)-F(x)}{h}\right)=2t.
\]
As $h\to0$, the quotients stay in the closed set
$2^{2t}\Ztwo^\times$, proving the derivative formula. At $x=0$ or
$-1$, the pair formula instead gives
\[
 v_2(F(x+h)-F(x))=3v_2(h)\qquad(v_2(h)\ge2).
\]
The corresponding difference quotients tend to zero. Thus these are
exactly the critical points. To identify the order at either point,
take the first nonzero coefficient in its local Taylor series. For
sufficiently large $v_2(h)$, that term controls the valuation of the
series. The displayed equality forces its degree to be $3$ and its
coefficient to be a unit.

We next determine the image. We use the elementary fact that every
isometry $g:\Ztwo\to\Ztwo$ is onto. For each $j\ge1$, the induced
map modulo $2^j$ is an injection of a finite set into itself, hence a
permutation. For any target, its nonempty inverse congruence classes
form nested compact subsets of $\Ztwo$; their intersection contains a
preimage.

For $r=1,2$, define
\[
 g_r(u)=\frac{F(4u+r)-F(r)}{16}\qquad(u\in\Ztwo).
\]
The local coefficient bound makes this integral. Both arguments in
each difference have $t$-value $1$, so \eqref{eq:padic-pairs} gives
\[
 v_2(g_r(u)-g_r(v))=v_2(u-v).
\]
Thus $g_r$ is onto. Since $F(1)=3$ and $F(2)=13$, the images of
$1+4\Ztwo$ and $2+4\Ztwo$ are exactly the two balls in
\eqref{eq:exact-image}.

For each $a\ge2$, parametrize the shell $v_2(x)=a$ by
$x=2^a(1+2u)$, and put
\[
 g_a(u)=\frac{F(2^a(1+2u))-F(2^a)}{2^{3a+1}}.
\]
Here \eqref{eq:padic-pairs} again makes $g_a$ an integral isometry of
$\Ztwo$ onto itself. Its shell image is therefore
\[
 F(2^a)+2^{3a+1}\Ztwo=1+2^{3a}\Ztwo^\times,
\]
where the last equality follows from
Proposition~\ref{prop:basepoints}. Reflection gives the shell image
near $-1$ as $-1+2^{3a}\Ztwo^\times$. Together with the two points
$0,-1$, these shells and the classes $1,2$ modulo $4$ partition
$\Ztwo$, proving \eqref{eq:exact-image}.

The two large image balls are $3$ and $13$ modulo $16$, whereas
the two shell families are $1$ and $15$ modulo $16$. Within either
family, different $a$ give different valuations relative to its
center; the centers themselves are excluded. Hence the union is
disjoint. Each shell image is a ball of measure $2^{-3a-1}$, and
singletons have measure zero. It follows that
\[
 m(F(\Ztwo))=\frac2{16}+2\sum_{a\ge2}2^{-3a-1}
 =\frac18+\frac1{56}=\frac17.
\]
\end{proof}
